\section{The Due Made Common}

\subsection{A received day}

St.~Justin Martyr's second-century account of the Sunday assembly holds together readings, common prayer, Eucharistic thanksgiving, Communion, and aid to those in need (\emph{First Apology} 67). Sunday is the Lord's Day: Christ rose on the first day of the week, and the Christian assembly receives that day as the first feast and sign of the new creation (CCC 2174). The reality later specified by law is already Paschal, sacramental, corporate, and merciful.

The juridical precept is understood from a reality it did not create. Early Christians experienced the Sunday gathering as an inner necessity before later ecclesiastical law gave the duty more explicit universal form (\emph{Dies Domini} 46--48). The Second Vatican Council calls Sunday the original feast day (\emph{Sacrosanctum Concilium} 106). Neither historical spontaneity nor later obligation exhausts the practice.

Because Eucharistic worship is the act of a body rather than an aggregate of private intentions, its common temporal form cannot be generated anew by each individual will. Sunday gathers persons who did not choose one another into one thanksgiving, interrupts labor's claim to ultimacy, and makes room for joy, rest, and mercy. A common minimum protects a good that private intentions alone cannot secure.

For members of the Latin Church whom the law binds, that minimum includes participation in Mass on Sunday or the preceding evening and abstention from work and affairs that impede worship, the day's proper joy, or suitable rest. The law itself recognizes impossibility, serious reason, and relief by competent authority; personal culpability further requires the knowledge and consent belonging to any grave sin. Reception of Communion at every Sunday Mass is not part of the Sunday precept.\endnote{Current universal Latin discipline as checked through 23 July 2026: CIC cc.~1, 11, 1245, 1246 \S1, 1247, and 1248; CCC 1857--1861 and 2180--2188. Canon 1248 \S1 permits fulfillment at Mass in any Catholic rite on the feast or preceding evening. Canon 1248 \S2 treats impossibility; canon 1245 identifies competent dispensers and commutation. Particular law, proper law, jurisdiction, individual facts, and Eastern Catholic law may change the analysis. The complete audit remains in the article's current-law matrix.}

The precision matters because justice does not authorize invented burdens. It also does not reduce Sunday to attendance. Works of mercy belong positively to the day; necessary care is not a profanation, and Christians should resist imposing avoidable labor upon those whose poverty deprives them of ordinary leisure (CCC 2185--2187). The precept marks a floor. It does not measure the fullness of worship or love.

\subsection{The ordinary means received}

The same structure appears more broadly in sacramental life. The sacraments are not techniques by which a creature gains leverage over God. They are Christ's instituted, ecclesial means by which grace is signified and given; for believers, they are the ordinary means of salvation in the differentiated manner proper to each sacrament (CCC 1127--1129). A Catholic cannot make private sincerity a substitute for the means Christ has given or treat ecclesial reception as an optional ornament to a self-designed spiritual life.

When justice has failed, the first truthful action is not to disguise the failure as devotion. Reconciliation names sin, receives sacramental absolution, and includes the purpose of amendment and satisfaction; it restores communion by Christ's mercy rather than by a performance that makes the sinner even with God (CCC 1422, 1440--1460). The point here is not to inventory sacramental discipline. It is that the received order includes a received way back.

The due made common is therefore not anonymous bureaucracy. It is a protection against the fantasy that the individual creature may define both the relation and its adequate response. The rite can be approached first because it is owed. Part III will ask how the very same reception can ripen from bare observance into hunger.
